\section{Additional Exercises}

\subsection{Lorentz Force Law}

\begin{exercise}\label{ex:spacetimematter-1}
Recall from the lecture that for a particle coupling to the electromagnetic potential, we have
\bse
    m\big(\nabla_{v_{\gamma}}v_{\gamma}\big)^a = q{F^a}_b{v_{\gamma}}^b,
\ese
where $v_{\gamma}$ is the velocity of a particle of mass $m$ and charge $q$.

Now "$1+3$"-decompose this equation in components with respect to the frame of an observer.
\end{exercise}
\hyperref[sol:spacetimematter-1]{\small \textit{Solution on p.~\pageref*{sol:spacetimematter-1}}}

\begin{exercise}\label{ex:spacetimematter-2}
Using the definitions $E_{\a}$ := $F_{\a0}$ for the electric field and $B^{\a} := \frac{1}{2}\varepsilon^{\a\rho\sig}F_{\rho\sig}$ for the magnetic 2
field seen by an observer, bring the right hand side of the above equation to the familiar form of the Lorentz force law for a particle of charge $q$ and spatial velocity
\bse
    \mathbf{v} := \bigg(\frac{e^{\a}:v_{\del}}{e^0:v_{\del}}\bigg) e_{\a} \qquad (\a=1,2,3\text{ and careful: the denominator was forgotten in lectures})
\ese
that the observer detects for the particle.

\textit{Hint: $(a\times b)^{\a} = g^{\a\mu}\varepsilon_{\mu\rho\sig}a^{\rho}b^{\sig}$, $\varepsilon_{123}=1$ and $\epsilon^{123}=1$.}
\end{exercise}
\hyperref[sol:spacetimematter-2]{\small \textit{Solution on p.~\pageref*{sol:spacetimematter-2}}}

\subsection{Which Curvature Can Feature In Einstein's Equations?}

\textit{This question asks to show that the differential Bianchi identity holds for the Riemann tensor. This was given as an exercise in the lecture.

If the reader couldn't work it out there, \href{https://math.stackexchange.com/questions/1494262/direct-proof-of-the-second-bianchi-identity}{this link} should prove helpful. Note the notation is slightly different in the link, but of course the answer is the same.

I didn't include this link in the notes to hopefully avoid temptation of just looking up the answer.}

\begin{exercise}\label{ex:spacetimematter-3}
The above component-free version can equivalently be written as
\bse
    {R^w}_{zab;c} + {R^w}_{zbc;a} + {R^w}_{zca;b} = 0.
\ese
Using this result, show that by appropriate contractions one obtains
\bse
    (\nabla_aG)^{ab}=0.
\ese
\end{exercise}
\hyperref[sol:spacetimematter-3]{\small \textit{Solution on p.~\pageref*{sol:spacetimematter-3}}}
